\documentclass[12pt,a4paper]{article}
\usepackage[brazil]{babel}
\usepackage[utf8]{inputenc}
\usepackage[T1]{fontenc}
\usepackage{lmodern}
\usepackage{geometry}
\geometry{top=3cm, bottom=2.5cm, left=3cm, right=2.5cm}
\usepackage{graphicx}
\usepackage{array}
\usepackage{tabularx}
\usepackage{threeparttable}
\usepackage{float}
\usepackage{caption}
\usepackage{hyperref}
\usepackage{listingsutf8}
\usepackage{xcolor}
\usepackage{enumitem}

% Configurações de caption e ambiente quadro
\captionsetup[table]{position=top, font=small}
\captionsetup[quadro]{justification=centering, labelsep=endash, width=0.75\linewidth, font=small}
\newfloat{quadro}{tbhp}{loq}
\floatname{quadro}{Quadro}

% Comando para fonte padronizada
\newcommand{\oautor}{Wagner Oliveira de Araujo}

% Hyperref
\hypersetup{colorlinks=true, linkcolor=blue, citecolor=blue, urlcolor=blue, filecolor=blue}

% Listings com suporte nativo a UTF-8 e mapeamento de acentos (resolve o erro do TeXnicCenter)
\lstset{
  language=Python,
  extendedchars=true,
  inputencoding=utf8,
  literate=
    {á}{{\'a}}1 {ã}{{\~a}}1 {ç}{{\c{c}}}1 {é}{{\'e}}1 {ê}{{\^e}}1
    {í}{{\'i}}1 {ó}{{\'o}}1 {õ}{{\~o}}1 {ú}{{\'u}}1
    {Á}{{\'A}}1 {Ã}{{\~A}}1 {Ç}{{\c{C}}}1 {É}{{\'E}}1 {Í}{{\'I}}1
    {Ó}{{\'O}}1 {Õ}{{\~O}}1 {Ú}{{\'U}}1,
  basicstyle=\ttfamily\small,
  breaklines=true,
  frame=single,
  backgroundcolor=\color{gray!10},
  keywordstyle=\color{blue},
  commentstyle=\color{gray},
  stringstyle=\color{red!70!black},
  numbers=left,
  numberstyle=\tiny\color{gray},
  stepnumber=1,
  numbersep=5pt,
  showstringspaces=false
}

\begin{document}

\title{\textbf{Projeto Final: Sistema de Regressão Espectral com Bagging e Subagging}\\
\small Reengenharia e Implementação em Streamlit/Python}
\author{Wagner Oliveira de Araujo}
\date{IFG --- 2026}
\maketitle
\thispagestyle{empty}

\newpage
\tableofcontents
\newpage

\section{Levantamento de Requisitos}

O levantamento visou catalogar, classificar e priorizar as funcionalidades do sistema legado (Vaadin/Java) para orientar a reengenharia em Python/Streamlit, garantindo rastreabilidade e aderência às necessidades de pesquisadores e alunos de quimiometria.

\subsection{Inventário e Classificação de Requisitos Funcionais}

\begin{quadro}[ht]
\begin{center}
\small
\caption{Inventário das funções do sistema}
\label{quadro:funcoes}
\begin{threeparttable}
\begin{tabularx}{\textwidth}{|l|l|X|} \hline
\textbf{Categoria} & \textbf{Funcionalidade} & \textbf{Descrição Técnica} \\ \hline
Cálculos & RLM Simples (Baseline) & Regressão Linear Múltipla padrão sem técnicas de ensemble \\ \hline
Cálculos & RLM + Bagging & Ensemble com amostragem \textit{com} reposição (bootstrap) \\ \hline
Cálculos & RLM + Subagging + CV & Ensemble com amostragem \textit{sem} reposição + otimização por validação cruzada 5-Fold \\ \hline
Configuração & Parâmetro \textit{m} (Bags) & Define número de amostras/ensembles para o método Bagging \\ \hline
Configuração & Parâmetro \textit{k} (Subagging) & Define fração ou quantidade para amostragem no Subagging \\ \hline
Visualização & Métricas de Desempenho & Exibição de R², RMSE, MAE calculados no conjunto de teste \\ \hline
Visualização & Gráficos Comparativos & \textit{Scatter plots} (previsto $\times$ real), resíduos e comparação entre modelos \\ \hline
Organização & Abas de Navegação & Separação lógica: ``Arquivos'', ``Configuração'', ``Resultados'', ``Gráficos'' \\ \hline
Exportação & \textit{Não identificada} & \textbf{Obs.:} sistema não apresenta botões de download de resultados \\ \hline
\end{tabularx}
\begin{tablenotes}[flushleft]
\footnotesize
\item Fonte:~\oautor
\end{tablenotes}
\end{threeparttable}
\end{center}
\end{quadro}

\begin{quadro}[ht]
\begin{center}
\small
\caption{Entradas necessárias para execução do sistema}
\label{quadro:entradas}
\begin{threeparttable}
\begin{tabularx}{\textwidth}{|l|l|X|} \hline
\textbf{Tipo} & \textbf{Nome/Parâmetro} & \textbf{Descrição} \\ \hline
Arquivo & \texttt{Xcal.txt} & Matriz de espectros de calibração (variáveis independentes) \\ \hline
Arquivo & \texttt{Ycal.txt} & Vetor de valores de referência para calibração \\ \hline
Arquivo & \texttt{Xteste.txt} & Matriz de espectros para predição/teste \\ \hline
Arquivo & \texttt{Yteste.txt} & Vetor de valores reais para validação do conjunto de teste \\ \hline
Parâmetro & \textit{m} (nº de bags) & Inteiro $\ge 1$. Usado apenas nos métodos de ensemble \\ \hline
Parâmetro & \textit{k} (Subagging) & Decimal $(0, 1]$. Usado apenas no método Subagging \\ \hline
Config. & Validação Cruzada & Fixo: 5-Fold. Parâmetro interno de otimização do modelo \\ \hline
\end{tabularx}
\begin{tablenotes}[flushleft]
\footnotesize
\item Fonte:~\oautor
\end{tablenotes}
\end{threeparttable}
\end{center}
\end{quadro}

\begin{quadro}[ht]
\begin{center}
\small
\caption{Saídas esperadas do sistema}
\label{quadro:saidas}
\begin{threeparttable}
\begin{tabularx}{\textwidth}{|l|l|X|} \hline
\textbf{Tipo} & \textbf{Conteúdo} & \textbf{Formato de Exibição} \\ \hline
Numérico & R², RMSE, MAE (teste) & Tabela ou cards destacados \\ \hline
Gráfico & Previsto vs. Real & Scatterplot com reta de ajuste \\ \hline
Gráfico & Resíduos da Predição & Gráfico de dispersão ou histograma \\ \hline
Comparativo & Desempenho entre modelos & Tabela comparativa ou gráfico de barras \\ \hline
Feedback & Mensagens de erro/validação & Alertas em destaque na interface \\ \hline
Status & Indicador de processamento & Spinner ou barra de progresso \\ \hline
\end{tabularx}
\begin{tablenotes}[flushleft]
\footnotesize
\item Fonte:~\oautor
\end{tablenotes}
\end{threeparttable}
\end{center}
\end{quadro}

\subsection{Classificação por Prioridade}

\begin{quadro}[ht]
\begin{center}
\small
\caption{Requisitos essenciais do sistema (Must Have)}
\label{quadro:essenciais}
\begin{threeparttable}
\begin{tabularx}{\textwidth}{|l|X|X|} \hline
\textbf{ID} & \textbf{Requisito} & \textbf{Justificativa} \\ \hline
E01 & Upload e validação dos 4 arquivos & Dados são pré-requisito para qualquer cálculo estatístico \\ \hline
E02 & Execução da RLM Simples (baseline) & Referência obrigatória para comparação de desempenho \\ \hline
E03 & Cálculo de métricas básicas (R², RMSE, MAE) & Mínimo necessário para avaliação quantitativa da qualidade \\ \hline
E04 & Validação dimensional (X vs Y) & Previne erros de execução e resultados inválidos \\ \hline
E05 & Exibição clara dos resultados & Permite decisão técnica fundamentada em evidências \\ \hline
\end{tabularx}
\begin{tablenotes}[flushleft]
\footnotesize
\item Fonte:~\oautor
\end{tablenotes}
\end{threeparttable}
\end{center}
\end{quadro}

\begin{quadro}[ht]
\begin{center}
\small
\caption{Requisitos desejáveis do sistema (Should Have)}
\label{quadro:desejaveis}
\begin{threeparttable}
\begin{tabularx}{\textwidth}{|l|X|X|} \hline
\textbf{ID} & \textbf{Requisito} & \textbf{Benefício} \\ \hline
D01 & Gráficos interativos (zoom, tooltip) & Facilita exploração visual e apresentação de resultados \\ \hline
D02 & Exportação de resultados (CSV, PDF, PNG) & Permite documentação e reprodutibilidade da análise \\ \hline
D03 & Tooltips explicativos para \textit{m}, \textit{k} & Reduz curva de aprendizado para usuários não especialistas \\ \hline
D04 & Pré-processamento básico (normalização) & Melhora robustez com dados espectrais reais \\ \hline
D05 & Histórico de execuções na sessão & Permite comparar configurações sem reenviar arquivos \\ \hline
\end{tabularx}
\begin{tablenotes}[flushleft]
\footnotesize
\item Fonte:~\oautor
\end{tablenotes}
\end{threeparttable}
\end{center}
\end{quadro}

\begin{quadro}[ht]
\begin{center}
\small
\caption{Requisitos extras para fases futuras (Could Have)}
\label{quadro:extras}
\begin{threeparttable}
\begin{tabularx}{\textwidth}{|l|X|X|} \hline
\textbf{ID} & \textbf{Requisito} & \textbf{Cenário de Uso} \\ \hline
X01 & Autenticação e projetos privados & Ambientes corporativos ou pesquisa com dados sensíveis \\ \hline
X02 & Integração com APIs externas & Pipelines automatizadas de quimiometria \\ \hline
X03 & Suporte a formatos (.csv, .xlsx, .mat) & Flexibilidade para ingestão de diferentes fontes \\ \hline
X04 & Otimização automática de hiperparâmetros & Reduz necessidade de conhecimento técnico avançado \\ \hline
X05 & Dashboard comparativo multi-rodada & Análise estatística de estabilidade dos ensembles \\ \hline
\end{tabularx}
\begin{tablenotes}[flushleft]
\footnotesize
\item Fonte:~\oautor
\end{tablenotes}
\end{threeparttable}
\end{center}
\end{quadro}

\subsection{Tratamento de Exceções e Requisitos Não Funcionais}

\begin{quadro}[ht]
\begin{center}
\small
\caption{Tratamento de exceções e usabilidade}
\label{quadro:excecoes}
\begin{threeparttable}
\begin{tabularx}{\textwidth}{|l|X|X|} \hline
\textbf{Cenário} & \textbf{Comportamento do Sistema} & \textbf{Mitigação Proposta} \\ \hline
Arquivo inválido & Alerta vermelho + destaque no campo & Validação dimensional \textit{fail-fast} antes do processamento \\ \hline
Parâmetros fora do intervalo & Validação em tempo real + mensagem & Sliders com valores recomendados e limites explícitos \\ \hline
Timeout de processamento & Mensagem orientativa de otimização & Barra de progresso com etapas descritivas \\ \hline
Navegação sem executar & Persistência de uploads na sessão & Estado mantido em \texttt{st.session\_state} \\ \hline
Resultados não exportáveis & Impede documentação externa & Botão ``Baixar Relatório'' com métricas + gráficos \\ \hline
\end{tabularx}
\begin{tablenotes}[flushleft]
\footnotesize
\item Fonte:~\oautor
\end{tablenotes}
\end{threeparttable}
\end{center}
\end{quadro}

\begin{quadro}[ht]
\begin{center}
\small
\caption{Especificação de requisitos não funcionais}
\label{quadro:nao-funcionais}
\begin{threeparttable}
\begin{tabularx}{\textwidth}{|l|X|X|} \hline
\textbf{Categoria} & \textbf{Especificação} & \textbf{Critério de Aceite} \\ \hline
Performance & Tempo de cálculo (RLM/Ensemble) & RLM $< 10\,$s; Ensemble $< 2\,$min (CV 5-Fold) \\ \hline
Performance & Renderização de gráficos & $< 3\,$s após conclusão do cálculo \\ \hline
Usabilidade & Acessibilidade (WCAG 2.1 AA) & Contraste, labels, navegação por teclado \\ \hline
Usabilidade & Clareza textual e recuperação de erro & Mensagens orientam correção, não apenas apontam falha \\ \hline
Segurança & Isolamento de sessão e sanitização & Dados não vazam; validação de extensão/conteúdo \\ \hline
Escalabilidade & Crescimento de usuários/dados & Container com load balancer; processamento em batch \\ \hline
\end{tabularx}
\begin{tablenotes}[flushleft]
\footnotesize
\item Fonte:~\oautor
\end{tablenotes}
\end{threeparttable}
\end{center}
\end{quadro}

\section{Projeto Lógico}

A reengenharia adotou o modelo clássico de desenvolvimento, promovendo a separação de responsabilidades e a eliminação de acoplamentos presentes no código legado.

\begin{quadro}[ht]
\begin{center}
\small
\caption{Diagnóstico do código legado e arquitetura proposta}
\label{quadro:diagnostico}
\begin{threeparttable}
\begin{tabularx}{\textwidth}{|l|X|X|} \hline
\textbf{Aspecto} & \textbf{Estado Atual (app.py)} & \textbf{Solução na Reengenharia} \\ \hline
Arquitetura & Monolítica (UI + lógica + parsing) & Camadas separadas: Apresentação, Negócio, Dados \\ \hline
Validação & \texttt{try/except} genérico & Função \texttt{validate\_inputs()} com regras explícitas \\ \hline
Performance & Laços \texttt{for} em CV, sem cache & \texttt{@st.cache\_data}, vetorização NumPy, seed exposto \\ \hline
Tipagem/Doc & Ausente em funções críticas & Type hints + docstrings padrão NumPy \\ \hline
Reprodutibilidade & \texttt{seed=42} hardcoded & Parâmetro configurável + registro em logs \\ \hline
\end{tabularx}
\begin{tablenotes}[flushleft]
\footnotesize
\item Fonte:~\oautor
\end{tablenotes}
\end{threeparttable}
\end{center}
\end{quadro}

O fluxo lógico segue: \texttt{Upload $\rightarrow$ Validação dimensional $\rightarrow$ Configuração de $m$ e $k$ $\rightarrow$ Execução com CV 5-Fold $\rightarrow$ Exibição de resultados e gráficos}. A camada de negócios foi isolada para garantir que a matemática da regressão (pseudoinversa, bootstrap, subamostragem) seja testável independentemente da interface.

\section{Projeto Físico}

\begin{quadro}[ht]
\begin{center}
\small
\caption{Estrutura modular e responsabilidades}
\label{quadro:estrutura}
\begin{threeparttable}
\begin{tabularx}{\textwidth}{|l|X|} \hline
\textbf{Módulo/Arquivo} & \textbf{Função Específica} \\ \hline
\texttt{app.py} & Camada de apresentação (apenas UI + roteamento de abas) \\ \hline
\texttt{src/data\_loader.py} & Parsing de \texttt{.txt}, validação dimensional, dataclasses \\ \hline
\texttt{src/models.py} & Implementação de RLM, Bagging e Subagging (vetorizado) \\ \hline
\texttt{src/metrics.py} & Cálculo de RMSE, MAE e R² com tipos explícitos \\ \hline
\texttt{src/optimizers.py} & Validação cruzada 5-Fold e busca de hiperparâmetros \\ \hline
\texttt{tests/} & Suíte de testes unitários (\texttt{pytest}) para cada módulo \\ \hline
\end{tabularx}
\begin{tablenotes}[flushleft]
\footnotesize
\item Fonte:~\oautor
\end{tablenotes}
\end{threeparttable}
\end{center}
\end{quadro}

O ambiente de execução utiliza Python 3.10+, Streamlit 1.30+, NumPy/Pandas para cálculo, Plotly para visualização interativa e \texttt{pytest} para validação automatizada. O cache de sessão e a alocação dinâmica de memória previnem \texttt{MemoryError} em datasets espectrais grandes.

\section{Implementação}

A implementação materializa o projeto lógico e físico em código Python, mantendo fidelidade matemática ao original e ampliando usabilidade e reprodutibilidade.

\subsection{Destaques da Implementação (\texttt{app.py})}
\begin{lstlisting}[language=Python]
def perform_bagging(Xcal, Ycal, Xteste, num_bags, with_replacement, fraction, seed=42):
    """RLM + Bagging / Subagging com vetorização e seed configurável."""
    n = Xcal.shape[0]
    sample_size = n if with_replacement else max(1, int(n * fraction))
    rng = np.random.default_rng(seed)
    Ypred = np.zeros(Xteste.shape[0])
    for _ in range(num_bags):
        idx = rng.integers(0, n, size=sample_size) if with_replacement else \
              rng.choice(n, size=sample_size, replace=False)
        beta = np.linalg.pinv(Xcal[idx]) @ Ycal[idx]
        Ypred += (Xteste @ beta) / num_bags
    return Ypred
\end{lstlisting}

\subsection{Estratégia de Testes e Validação}

\begin{quadro}[ht]
\begin{center}
\small
\caption{Estratégia de testes e critérios de aceite}
\label{quadro:testes}
\begin{threeparttable}
\begin{tabularx}{\textwidth}{|l|X|X|} \hline
\textbf{Tipo de Teste} & \textbf{Escopo} & \textbf{Meta de Cobertura} \\ \hline
Unitário & Funções isoladas (parse, metrics) & $\ge 80\%$ (\texttt{pytest-cov}) \\ \hline
Integração & Fluxo completo: Upload $\rightarrow$ CV $\rightarrow$ Gráficos & 100\% dos caminhos críticos \\ \hline
Performance & Tempo de resposta para datasets típicos & RLM $< 10\,$s; Ensemble $< 2\,$min \\ \hline
Aceite (UAT) & Validação por pesquisadores/alunos & SUS $\ge 75$ (Bom/Excelente) \\ \hline
\end{tabularx}
\begin{tablenotes}[flushleft]
\footnotesize
\item Fonte:~\oautor
\end{tablenotes}
\end{threeparttable}
\end{center}
\end{quadro}

\subsection{Rastreabilidade de Requisitos}

\begin{quadro}[ht]
\begin{center}
\small
\caption{Matriz de rastreabilidade de requisitos (RTM)}
\label{quadro:rtm}
\begin{threeparttable}
\begin{tabularx}{\textwidth}{|l|l|X|} \hline
\textbf{ID Requisito} & \textbf{Fase Clássica} & \textbf{Artefato de Validação} \\ \hline
RF01.1--01.5 & 1, 3, 5 & \texttt{data\_loader.py}, testes de upload \\ \hline
RF02.1--02.3 & 1, 2, 4 & \texttt{optimizers.py}, logs de execução \\ \hline
RF03.1--03.3 & 2, 3, 4 & \texttt{models.py}, testes unitários de métricas \\ \hline
RF04.1--04.3 & 2, 4, 5 & \texttt{metrics.py}, UAT visual (Plotly) \\ \hline
RF05.1--05.3 & 1, 3, 5 & UI/UX, feedback de erros, \texttt{st.progress} \\ \hline
\end{tabularx}
\begin{tablenotes}[flushleft]
\footnotesize
\item Fonte:~\oautor
\end{tablenotes}
\end{threeparttable}
\end{center}
\end{quadro}

A implementação resultante atende aos requisitos funcionais e não funcionais levantados, reduz a curva de aprendizado para usuários não especialistas e estabelece base técnica sólida para evoluções futuras (autenticação, pipelines automatizados, suporte a novos formatos).

\end{document}